\renewcommand{\S}{\mathbf S}

\section{Appendix C: Estimates in hyperbolic geometry}

In this appendix we will describe how to measure the geometry of a surface
 using transverse geodesic arcs.
We will then show how to compute the lengths of peripheral closed geodesics of the surface
 using these measurements.

Suppose $T$ is a compact hyperbolic surface with geodesic boundary.
The following lemma appears as the Corollary to section 3.3 of \cite{Ab}:

\begin{lem} \label{margulis}
There is an $\ezero > 0$ such that any two distinct closed geodesics on $T$ of length at most $\ezero$
 are simple and disjoint.
\end{lem}

Let $\S$ be a compact hyperbolic surface with geodesic boundary.
Then a \emph{transverse geodesic arc} for $\S$
 is a proper path of minimal length in its proper homotopy class.
If $\bfalpha$ is a path on $\S$,
it is a transverse geodesic arc if and only if
 it is a geodesic arc that meets $\partial \S$ orthogonally,
or,
equivalently,
the double of $\bfalpha \cup \overline \bfalpha$
 in $\S \cup \overline \S$
is a closed geodesic.

Let $S$ be a compact Riemann surface with boundary,
and endow $\Int S$ with its Poincar\'e metric.
The peripheral geodesics on $\Int S$ bound a compact surface
$\S$ with geodesic boundary,
called the \emph{convex core} of $S$.
There is a homeomorphism $h\from \S \to S$
that is isotopic through embeddings to the inclusion $\S \subset S$.
We can then form a weighted arc diagram $M^S$ on $S$ as follows:
 for $\alpha \in \AA(S)$,
we find the transverse geodesic arc $\bfalpha$ for $\S$
such that $h(\bfalpha) \sim \alpha$.
Then we let $M^S(\alpha) = - \log L(\bfalpha)$ if
 $L(\bfalpha) < \ezero/2$,
and $M^S(\alpha) = 0$ otherwise.
Then $M^S$ is supported on a set of disjoint arcs,
so $M^S$ is a weighted arc-diagram for $S$.
We let
$\M^S = \bigcup\{ \bfalpha : L(\bfalpha) < \ezero/2 \}$,
so $\M^S \subset  \S$ is the union of the short transverse geodesic arcs of $\S$.

We will call a non-peripheral simple closed geodesic
 a \emph{dividing} geodesic.
The following result appears in \cite{Ab}:

\begin{lem} \label{abikoff}
Let $T$ be a compact hyperbolic Riemann surface with bounded-length geodesic boundary.
Then either $T$ is a pair of pants,
or there is a bounded-length dividing geodesic on $T$.
\end{lem}

We say that a hyperbolic surface $T$ is \emph{symmetric}
 if it admits an isometric orientation-reversing involution,
which we will denote by complex conjugation: ``$z \mapsto \overline z$''.
Then we let $E_T = \{ z \in T : z = \overline z \}$,
 and $E_T$ will be a union of (simple) closed geodesics and transverse geodesic arcs.
($E_T$ depends implicitly on on the choice of involution.
Whenever we say ``symmetric hyperbolic surface'' we will mean
 ``symmetric hyperbolic surface and choice of involution.'')
Note that $T \sm E_T$ has two components,
call them $A_T$ and $\overline{A_T}$,
which are mapped to each other by $z \mapsto \overline z$.
We prove a symmetric version of Lemma \ref{abikoff}:

\begin{lem} \label{symmetric abikoff}
Every symmetric compact hyperbolic surface $T$ with \break bounded-length geodesic boundary
 has a bounded-length symmetric pair of pants decomposition.
\end{lem}
\begin{proof}
It suffices to find a single bounded-length symmetric dividing geodesic on the surface,
or a symmetric pair of disjoint bounded-length dividing geodesics,
because then we can cut the surface $T$ along that geodesic or pair of geodesics,
and repeat.

By Theorem \ref{abikoff},
unless $T$ is a pair of pants,
there is a dividing geodesic $\gamma$ of bounded length on $T$.
If $\gamma \cap E_T = \emptyset$,
then $\gamma \cap \overline \gamma = \emptyset$,
and we are done.
Likewise, if $\gamma \subset E_T$,
then $\gamma$ is symmetric,
and we are done.
Otherwise,
let $\eta$ be a component of $\gamma \cap \Cl A_T$.
Then $\eta \cup \overline \eta$
is a non-trivial non-peripheral simple closed curve
so we let $\tau$ be the dividing
 geodesic homotopic to $\eta$.
Then $L(\tau) \le 2 L(\eta) < 2 L(\gamma)$,
so $\tau$ is the desired object.
\end{proof}

We now prove two basic estimates on transverse geodesic arcs on pairs of pants.
We denote by
$[x, y, z, r, s, t]$
the right-angled hyperbolic hexagon
with lengths $x, y, z, r, s, t$ in that order.
We will omit lengths that are not specified,
so for example $[a,,b,,c,]$
denotes the right-angled hexagon with alternating side lengths $a$, $b$, and $c$.
We first estimate the length of one side in a hyperbolic right-angled hexagon,
in terms of the lengths of three alternating sides:
\begin{lem}\label{fenchel}
Let $[a, c', b, , c, ,]$ be a hyperbolic right-angled hexagon,
and suppose that $a, b, c \le r$.
Then $c' = -\log a - \log b + O(1;r)$.\footnote{Notation $O(1;r)$ stands for a quantity  bounded
     in terms of  $r$. }
\end{lem}
\begin{proof}
We use the formula (from \cite{Fenchel:Elementary}):
\begin{equation}\label{hexagon}
\cosh c' = \frac{\cosh c + \cosh a \cosh b}{\sinh a \sinh b}
         = e^{O(1; r)} \frac 1 {ab}
\end{equation}
and recall that $\cosh^{-1} x = \log x + O(1; t)$
whenever $x \ge t > 0$.
\end{proof}
We let $\PP(a, b, c)$ denote the pair of pants with cuff lengths $a$, $b$, and $c$.
\begin{lemma}
We have the following two estimates:
\begin{enumerate}
\item
If $P = \PP(a, b, c)$ is a pair of pants,
and $\gamma$ is the transverse geodesic arc that connects $a$ and $b$,
then
$$
|\gamma| = -\log a - \log b + O(1),
$$
for $a, b, c \le C_0$.
\item
If $\gamma$ is the transverse geodesic arc that connects $a$ and $a$ in $P(a, b, b)$,
then
$$
|\gamma| = - 2 \log a + O(1)
$$
for $a, b \le C_0$.
\end{enumerate}
\end{lemma}
\begin{proof}
We prove each of the above:
\begin{enumerate}
\item
We cut $P$ along the three pairwise transversals
into two right-angled hexagons.
By formula (\ref{hexagon}) these hexagons are equal,
and hence each has type
$[a/2, \gamma,b/2, \, , c/2, \,]$.
Apply now Lemma \ref{fenchel}.

\item
We let $\eta$ be the transversal from the length $a$ cuff to one other.
Then we cut along $\eta$ and $\gamma$ to obtain the right-angled hexagon
$[a/4, \gamma, a/4, \eta, b, \eta]$,
and then apply Lemma \ref{fenchel}. \qedhere
\end{enumerate}
\end{proof}

Given a closed geodesic $\gamma$ and an arc $\alpha\in \AAA$,
let
$\pairing{ \gamma, \alpha }$
stand for the intersection number of $\gamma$ with $\alpha$,
i.e.,
the minimal number of intersections of $\gamma$ with the paths representing $\alpha$.
Given a weighted arc diagram $W=\sum W(\alpha)\alpha$,
we can define the intersection number
$\pairing{ \gamma, W }=\sum W(\alpha)\pairing{ \gamma, \alpha }$
by linearity.
We can now prove the main theorem of this appendix:
\begin{thm} \label{peripheral lengths}
Let $S$ be a compact Riemann surface with boundary,
and endow $\Int S$ with its Poincar\'e metric.
Suppose that $\gamma$ is a peripheral closed geodesic for $S$.
Then
$$
L(\gamma) = 2 \left< M^S, \gamma \right> + O(1; \chi(S)).
$$
\end{thm}

\begin{proof}
We let $\S$ be the convex core of $S$.
We find a symmetric bounded-length pair of pants decomposition for $\S \cup \overline \S$
extending $\M^S\cup \bar \M^S$.
Then we can write $\gamma = \bigcup t_i$,
where the segments $t_i$ are interior-disjoint,
and each is a transverse arc of one of the pairs of pants.
Then $L(t_i) = -\log a_i - \log b_i + O(1)$,
where $a_i$ and $b_i$ are the lengths of the cuffs that $\gamma$ connects
(possibly the same cuff).
Therefore
$$
L(\gamma) = 2 \sum - \log a_i + O(1)
$$
where the $a_i$ are the lengths of the cuffs that $\gamma$ crosses,
counted with multiplicity.
But
\begin{equation}
2 \sum -\log a_i = 2 \left< M^S, \gamma \right> + O(1). \qedhere
\end{equation}
\end{proof}

